একটি \emph{টুরিং যন্ত্র} হলো গণনা করার এক ধরনের আদর্শায়িত যন্ত্র।
এতে একমুখী অসীম একটি \emph{ফিতা} থাকে, যা বর্গাকার ঘরে বিভক্ত।
প্রতিটি ঘরে আগে থেকে স্থির করা একটি বর্ণমালার একটি \emph{প্রতীক}
থাকতে পারে। যন্ত্রটি ফিতা বরাবর একটি \emph{পাঠ-লেখ হেড}
সরিয়ে কাজ করে। যন্ত্রটি আগে থেকে নির্দিষ্ট কয়েকটি \emph{দশার}
যেকোনো একটিতে থাকতে পারে। পাঠ-লেখ হেডের কাজগুলি একটি নির্দেশ-সেট
দিয়ে নির্ধারিত হয়। প্রতিটি নির্দেশ প্রযোজ্য হয় যন্ত্রটি একটি
নির্দিষ্ট দশায় থাকলে এবং একটি নির্দিষ্ট প্রতীক পড়লে। নির্দেশটি বলে
দেয় যন্ত্রটি বর্তমান ঘরে কোন প্রতীক লিখবে, পাঠ-লেখ হেডটিকে এক ঘর
বাঁয়ে সরাবে, ডাইনে সরাবে, না একই ঘরে রাখবে, এবং কোন দশায়
যাবে। ফিতায় প্রতীকের অনুক্রম হিসেবে কোনো \emph{নিবেশ} লেখা থাকলে
এবং পাঠ-লেখ হেডটি নিবেশের বাঁদিকের প্রথম ঘরটি পড়ছে এমন অবস্থায়
যন্ত্রটিকে নির্দিষ্ট আরম্ভিক দশায় বসালে, নির্দেশ-সেটটি ধাপে ধাপে
যন্ত্রের \emph{কনফিগারেশনের} একটি অনুক্রম নির্ধারণ করবে। একটি
কনফিগারেশনে রয়েছে ফিতায় যা লেখা আছে, পাঠ-লেখ হেডের অবস্থান এবং
যন্ত্রের দশা। যন্ত্রটি এমন কোনো কনফিগারেশনে পৌঁছালে যেখানে পঠিত
প্রতীক ও বর্তমান দশার ওই যুগলের জন্য নির্দেশ-সেটে কোনো নির্দেশ
নেই, যন্ত্রটি \emph{থামে}, এবং ফিতায় যা লেখা থাকে, সেটিই
\emph{নির্গম}।

টুরিং যন্ত্রের ফিতায় দাগের অনুক্রম দিয়ে খুব সহজেই সংখ্যা
উপস্থাপন করা যায়। একটি অপেক্ষক $\Nat \to \Nat$-কে
\emph{টুরিং-গণনসাধ্য} বলি যদি এমন একটি টুরিং যন্ত্র থাকে যা
নিবেশ হিসেবে~$n$-এর এককীয় উপস্থাপনা দিয়ে চালু করা হলেই
একসময় থামে এবং নির্গম হিসেবে ফিতায়~$f(n)$-এর এককীয়
উপস্থাপনা রেখে যায়। পরিচিত অনেক পাটিগাণিতিক অপেক্ষককে সহজেই
(বা কিছুটা কষ্টে) টুরিং-গণনসাধ্য দেখানো যায়। টুরিং যন্ত্র
ছাড়াও গণনার আরও অনেক মডেল প্রস্তাব করা হয়েছে। প্রত্যেক
ক্ষেত্রেই দেখা গেছে, ওই মডেলে যে পাটিগাণিতিক অপেক্ষকগুলি গণনা করা
যায়, সেগুলি টুরিং-গণনসাধ্যও। এই ফলকে \emph{চার্চ--টুরিং
থিসিসের} সমর্থন হিসেবে দেখা হয়। থিসিসটি বলে, কার্যকরভাবে গণনা
করা যায় এমন প্রতিটি পাটিগাণিতিক অপেক্ষকই টুরিং-গণনসাধ্য।
